% 结论章时间线图 — 第7章用
% 三阶段前瞻判断时间线：原型验证 → 有限部署 → 条件性扩张
% 明确为"前瞻判断时间线"，非历史事件回顾
\begin{tikzpicture}[
  phase/.style={rectangle, draw=gray!70, rounded corners, minimum width=3.0cm,
                minimum height=1.2cm, align=center, font=\small},
  judge/.style={text width=2.8cm, align=center, font=\footnotesize, text=black!70},
  arrow/.style={thick, draw=gray!60, ->, >=stealth},
  note/.style={font=\small\itshape, text=gray!50!black, align=center}
]
  % ===== 底部时间轴 =====
  \draw[thick, ->, >=stealth] (-0.3,0) -- (10,0);

  % 时间刻度
  \foreach \x/\l in {0/2026, 2.8/2028, 5.6/2030, 8.4/2035}
    \draw (\x,0.08) -- (\x,-0.08) node[below, font=\footnotesize] {\l};

  % ===== 阶段一：2026--2028 原型验证期 =====
  \node[phase, fill=gray!20] (p1) at (1.4,0.9) {\textbf{2026--2028}\\[2pt]\textbf{原型验证期}};
  \node[judge] at (1.4,-0.9) {经济性不可行\\技术验证有意义};

  % ===== 阶段二：2028--2030 有限部署期 =====
  \node[phase, fill=cyan!18!gray!14] (p2) at (4.2,0.9) {\textbf{2028--2030}\\[2pt]\textbf{有限部署期}};
  \node[judge] at (4.2,-0.9) {利基场景可行\\规模化不可行};

  % ===== 阶段三：2030--2035 条件性扩张期 =====
  \node[phase, fill=blue!16!gray!20] (p3) at (7.0,0.9) {\textbf{2030--2035}\\[2pt]\textbf{条件性扩张期}};
  \node[judge] at (7.0,-0.9) {取决于多参数\\同步改善概率};

  % ===== 阶段间演进箭头 =====
  \draw[arrow] (p1.east) -- (p2.west);
  \draw[arrow] (p2.east) -- (p3.west);

  % ===== 底部说明 =====
  \node[note] at (4.6,-2.0) {本图为前瞻判断时间线，基于第4--6章分析结论，不代表实际发生预测};
\end{tikzpicture}
